\documentclass[11pt,a4paper]{article}
\usepackage[utf8]{inputenc}
\usepackage[spanish]{babel}
\usepackage{amsmath, amssymb, amsfonts}
\usepackage{geometry}
\usepackage{xcolor}
\usepackage{booktabs}
\usepackage{tcolorbox}

\geometry{margin=2.5cm}

\title{\textbf{01. Preparación de Datos y Estimación de Parámetros (Etapa 1)}}
\author{Registro de Definiciones del Proyecto - Tesis DRL Portafolios}
\date{\today}

\begin{document}

\maketitle

Este documento detalla la ejecución técnica de la \textbf{Etapa 1} definida en la Memoria de Referencia (Documento 02), adaptándola a los activos SPX y CMC200 para un horizonte de $T=163$ semanas.

\section{Datos Fuente y Horizonte Temporal}

El proyecto utiliza series históricas semanales para $T=163$ periodos. Los insumos base son:
\begin{itemize}
    \item \textbf{Retornos Netos ($R_{i,t}$):} Observaciones semanales de SPX y CMC200.
    \item \textbf{Probabilidades HMM ($P_{i,k,t}$):} Señales exógenas de régimen (Bear/Bull).
\end{itemize}

\section{Partición del Dataset (Rigor DRL)}

Siguiendo la metodología \textit{walk-forward} de la Biblia, dividimos las 163 semanas para garantizar la capacidad de generalización:

\begin{table}[h!]
\centering
\begin{tabular}{@{}lll@{}}
\toprule
\textbf{Conjunto} & \textbf{Rango (semanas)} & \textbf{Propósito} \\ \midrule
Entrenamiento (Train) & 1 - 114 (70\%) & Aprendizaje de la política. \\
Validación (Val) & 115 - 138 (15\%) & Ajuste de hiperparámetros. \\
Prueba (Test) & 139 - 163 (15\%) & Evaluación out-of-sample vs GAMS. \\ \bottomrule
\end{tabular}
\end{table}

\section{Estimación de Parámetros Mix}

Para cada periodo $t$, se calculan las expectativas de mercado mediante una ventana expansiva ($1 \dots t$).

\subsection{Estimación de Momentos por Régimen}
Para cada activo $i$ bajo el régimen $k$, los momentos se estiman ponderando cada observación histórica por su probabilidad de pertenencia al régimen:

\begin{equation}
    \hat{\mu}_{i,k,t} = \frac{\sum_{m=1}^t P_{i,k,m} R_{i,m}}{\sum_{m=1}^t P_{i,k,m}}
\end{equation}

\begin{equation}
    \hat{\sigma}_{i,j,k,t} = \frac{\sum_{m=1}^t P_{i,k,m} (R_{i,m} - \hat{\mu}_{i,k,t})(R_{j,m} - \hat{\mu}_{j,k,t})}{\sum_{m=1}^t P_{i,k,m}}
\end{equation}

\subsection{Construcción de Variables Mix}
Integramos la incertidumbre del régimen en una única señal de media y riesgo:
\begin{equation}
    \mu_{mix, i,t} = \sum_{k} P_{i,k,t} \cdot \hat{\mu}_{i,k,t}
\end{equation}
\begin{equation}
    \sigma_{mix, i,j,t} = \sum_{k} P_{i,k,t} \cdot \hat{\sigma}_{i,j,k,t}
\end{equation}

Estas variables simplifican la tarea del agente al entregarle la información ya integrada.

\section{Normalización Z-Score de Observaciones}

Aplicamos estandarización expansiva para asegurar estabilidad numérica. Para cada variable $x_{mix, t}$, calculamos su versión normalizada:
\begin{equation}
    z(x)_{mix, t} = \frac{x_{mix, t} - \mu(x)_{mix, t}}{\sigma(x)_{mix, t}}
\end{equation}

\section{Parámetros del Modelo y Configuración Base}

Para asegurar la comparabilidad con el modelo GAMS original, se registran las siguientes constantes:

\begin{itemize}
    \item \textbf{Capital Inicial ($x_0$):} 10,000 USD.
    \item \textbf{Pesos Iniciales ($w_0$):} 50\% SPX / 50\% CMC200.
    \item \textbf{Costos de Transacción ($c_i$):} SPX ($0.5\%$), CMC200 ($1.0\%$).
    \item \textbf{Aversión al Riesgo ($\lambda$):}
    \begin{itemize}
        \item Grilla GAMS: $\{0.05, 0.10, 0.25, 0.50, 1.00\}$.
        \item \textbf{Selección DRL:} $\lambda = 0.10$.
    \end{itemize}
\end{itemize}

\end{document}
